\chapter{النمذجة الديناميكية (\LR{Dynamic Modeling})}
\label{ch:dynamics}

بعد استنتاج النموذج الحركي في الفصل السابق والذي يصف الحركة هندسياً، ينتقل التحليل في هذا الفصل إلى المستوى الديناميكي؛ حيث تتم دراسة القوى والعزوم المسببة لهذه الحركة. تعتبر النمذجة الديناميكية الدقيقة حجر الأساس في هندسة التحكم الآلي المتقدم، وبدونها لا يمكن تصميم خوارزميات تحكم قوية تعتمد على النموذج (\LR{Model-Based Control}) لمواجهة التداخلات الديناميكية وتغيرات الحمل.

سيتم في هذا الفصل اشتقاق النموذج الديناميكي للروبوت (\LR{SCARA}) بالاعتماد على صياغة لاغرانج-أويلر (\LR{Euler-Lagrange}). ولضمان دقة فيزيائية عالية، سيتم أخذ تغيرات الحمل بعين الاعتبار، بالإضافة إلى دمج ديناميكية المحركات الكهربائية للحصول على نموذج متكامل يمثل النظام الفعلي بدقة.


\section{النهج المعتمد للنمذجة الديناميكية والفرضيات الأساسية}
\label{sec:dynamic_approach}

يعتبر بناء نموذج ديناميكي دقيق تحدياً رياضياً، حيث تتشابك المعادلات بشكل كبير مع زيادة درجات الحرية. لتبسيط الحسابات الرياضية وتخفيف العبء الحسابي على وحدة التحكم دون الإخلال بالتمثيل الفيزيائي الدقيق للنظام، تم الاستناد إلى مجموعة من الفرضيات الهندسية والفيزيائية عند بناء النموذج الديناميكي لروبوت (\LR{SCARA}) المدروس:

\begin{enumerate}
    \item \textbf{تغير كتلة الوصلة الثالثة (تأثير الحمل):}
    بما أن المهمة الأساسية للروبوت هي نقل الحمولات (\LR{Pick and Place})، فإن الكتلة الفعالة للوصلة الثالثة لا تبقى ثابتة. لذلك، تم دمج كتلة الحمل $(m_L)$ بشكل صريح مع كتلة الوصلة الثالثة الأساسية $(m_3)$ لتشكيل كتلة مركبة جديدة متغيرة تعطى بالعلاقة $(m_3^* = m_3 + m_L)$. هذا الدمج يتطلب إعادة حساب شعاع مركز الثقل $(r_{p3})$ ومصفوفة القصور الذاتي للوصلة الثالثة في كل لحظة يتغير فيها الحمل.
    
    \item \textbf{المرجعية الصفرية للجاذبية الأرضية:}
    تم اعتبار الجملة الإحداثية المرتبطة بالمفصل الأول  هي الجملة المرجعية (\LR{Reference Frame}) لحسابات الطاقة الكامنة. وبناءً عليه، فإن شعاع الجاذبية الأرضية يؤثر عمودياً نحو الأسفل وبشكل موازٍ للمحور $Z_0$، ويُعطى بالصيغة الشعاعية $g = [0, 0, -g]^T$.
    
    \item \textbf{إهمال قوى كوريوليس والقوى النابذة:}
    في العديد من التطبيقات الصناعية لروبوتات (\LR{SCARA})، وبناءً على مسارات التشغيل التي لا تتطلب سرعات زاوية عالية، يمكن افتراض أن النظام يعمل ضمن سرعات حركية منخفضة إلى متوسطة. يتيح هذا الافتراض إهمال مصفوفة قوى كوريوليس والطرد المركزي $C(q, \dot{q})$ من النموذج الديناميكي العام، مما يبسط النموذج بشكل كبير.
\end{enumerate}

بناءً على صياغة لاغرانج-أويلر والفرضيات المذكورة أعلاه، يتبسط النموذج الديناميكي العام للروبوت (المجرد من المحركات) من شكله القياسي المعقد إلى الصيغة المصفوفية المخفضة التالية:

\begin{equation}
M(q)\ddot{q} + G(q) = \tau_{robot}
\label{eq:simplified_dynamic_model}
\end{equation}

حيث:
\begin{itemize}
    \item $M(q)$: مصفوفة العطالة أو القصور الذاتي للنظام بأبعاد $3 \times 3$.
    \item $G(q)$: شعاع قوى الجاذبية الأرضية بأبعاد $3 \times 1$.
    \item $\ddot{q}$: شعاع التسارعات المفصلية بأبعاد $3 \times 1$.
    \item $\tau_{robot}$: شعاع العزوم والقوى المحركة الفعالة عند المفاصل بأبعاد $3 \times 1$.
\end{itemize}

في الأقسام التالية، سيتم تفصيل الخطوات الرياضية لحساب مصفوفة العطالة وشعاع الجاذبية بالاعتماد على الثوابت الميكانيكية للوصلات واستخدام نظريات انتقال المحاور وعزوم القصور الذاتي.
\section{تحديد الثوابت الميكانيكية ومصفوفات القصور الذاتي للوصلات}
\label{sec:mass_inertia_parameters}

يتطلب حساب النموذج الديناميكي معرفة دقيقة بالثوابت الفيزيائية والميكانيكية لكل وصلة من وصلات الروبوت. يتضمن ذلك أطوال الوصلات، وكتلها، وتحديد متجهات مراكز الثقل، بالإضافة إلى حساب مصفوفات القصور الذاتي المتكاملة (\LR{Pseudo-Inertia Matrices}) التي تُشكل النواة الرياضية في صياغة لاغرانج-أويلر.

\subsection{المعايير الهندسية والكتلية للروبوت}
بناءً على التصميم الميكانيكي الفعلي لروبوت (\LR{SCARA}) المدروس، تم تحديد قيم الأطوال والكتل الثابتة للوصلات الثلاث، بالإضافة إلى دمج كتلة الحمل المتغير $(m_L)$ عند نهاية الوصلة الثالثة كما يلي:
\begin{itemize}
    \item أطوال الوصلات: $L_1 = 0.45 \text{ m}$، و $L_2 = 0.35 \text{ m}$، و $L_3 = 0.1 \text{ m}$.
    \item كتل الوصلات الأصلية (بدون حمل): $m_1 = 14 \text{ kg}$، و $m_2 = 18 \text{ kg}$، و $m_3 = 2 \text{ kg}$.
    \item الكتلة المركبة الإجمالية للوصلة الثالثة مع الحمل: $m_3^* = m_3 + m_L$.
\end{itemize}

\subsection{متجهات مراكز الثقل (\LR{Center of Mass Vectors})}
بفرض أن الوصلتين الأولى والثانية عبارة عن قضبان متجانسة الحجم والكثافة، فإن مركز ثقلهما الهندسي يقع في المنتصف تماماً على امتداد المحور الطولي الفعال. أما الوصلة الثالثة، فنظراً لكونها الجزء المسؤول عن حمل الأوزان المتغيرة، فإن مركز ثقلها الكلي ينزاح لحظياً بناءً على قيمة الحمل المرفوع $(m_L)$. 

باستخدام قانون العزوم السكونية لمركز الكتل المشترك :
\begin{equation}
R_c = \frac{\sum m_i r_i}{\sum m_i}
\end{equation}

تم تحديد متجهات مراكز الثقل $r_{p1}$ و $r_{p2}$ و $r_{p3}^*$ لكل وصلة منسوبة إلى جملة إحداثياتها المحلية كما يلي:

\begin{equation}
r_{p1} = \begin{bmatrix} \frac{L_1}{2} \\ 0 \\ 0 \end{bmatrix} , \quad 
r_{p2} = \begin{bmatrix} \frac{L_2}{2} \\ 0 \\ 0 \end{bmatrix} , \quad 
r_{p3}^* = \begin{bmatrix} 0 \\ 0 \\ z_{c3}^* \end{bmatrix}
\label{eq:com_vectors}
\end{equation}

حيث يمثل الإحداثي $z_{c3}^*$ موضع مركز الثقل المشترك للوصلة الثالثة مع الحمل على المحور العمودي، ويُحسب لحظياً عبر العلاقة التالية:
\begin{equation}
z_{c3}^* = \frac{m_3 \left(\frac{L_3}{2}\right) + m_L L_3}{m_3 + m_L}
\label{eq:zc3_star}
\end{equation}

\subsection{مصفوفات القصور الذاتي المتكاملة للوصلات (\LR{Pseudo-Inertia Tensors})}
في صياغة لاغرانج-أويلر، يتم دمج مصفوفة عزم القصور الذاتي الدوراني القياسية $I_i$ مع متجهات مركز الثقل والكتلة ضمن مصفوفة واحدة ممتدة بأبعاد $4 \times 4$ تُعرف بالمصفوفة المتكاملة $J_i$. تُعرّف هذه المصفوفة عموماً بالصيغة التالية:

\begin{equation}
J_i = \begin{bmatrix} 
\int x^2 dm & \int xy dm & \int xz dm & \int x dm \\ 
\int xy dm & \int y^2 dm & \int yz dm & \int y dm \\ 
\int xz dm & \int yz dm & \int z^2 dm & \int z dm \\ 
\int x dm & \int y dm & \int z dm & \int dm 
\end{bmatrix}
\end{equation}

بتطبيق هذا المفهوم وحساب التكاملات الحجمية للوصلات بناءً على توزيع الكتل المحدد ومراكز الثقل في المعادلات (\ref{eq:com_vectors})، تم استنتاج مصفوفات القصور الذاتي المتكاملة الجاهزة للتعويض المباشر في معادلات الحركة:

\textbf{1. مصفوفة القصور الذاتي للوصلة الأولى ($J_1$):}
\begin{equation}
J_1 = \begin{bmatrix}
\frac{1}{3}m_1 L_1^2 & 0 & 0 & \frac{L_1}{2}m_1 \\
0 & 0 & 0 & 0 \\
0 & 0 & 0 & 0 \\
\frac{L_1}{2}m_1 & 0 & 0 & m_1
\end{bmatrix}
\end{equation}

\textbf{2. مصفوفة القصور الذاتي للوصلة الثانية ($J_2$):}
\begin{equation}
J_2 = \begin{bmatrix}
\frac{1}{3}m_2 L_2^2 & 0 & 0 & \frac{L_2}{2}m_2 \\
0 & 0 & 0 & 0 \\
0 & 0 & 0 & 0 \\
\frac{L_2}{2}m_2 & 0 & 0 & m_2
\end{bmatrix}
\end{equation}

\textbf{3. مصفوفة القصور الذاتي المركبة للوصلة الثالثة مع الحمل ($J_3$):}
نظراً لتغير مركز الثقل، تم استخدام نظرية المحاور المتوازية لشتاينر (\LR{Steiner's Theorem}) لحساب انزياح عزم القصور الذاتي للاسطوانة والحمل عن المركز الجديد، حيث تعطى بالصيغة العامة $I_{new} = I_{center} + md^2$. بناءً على ذلك، تم صياغة المصفوفة المتكاملة المحدثة للحمل والوصلة الثالثة ككتلة واحدة بالشكل التالي:

\begin{equation}
J_3 = \begin{bmatrix}
0 & 0 & 0 & 0 \\
0 & 0 & 0 & 0 \\
0 & 0 & J_{3z} & -\left(m_3 \frac{L_3}{2} + m_L L_3\right) \\
0 & 0 & -\left(m_3 \frac{L_3}{2} + m_L L_3\right) & m_3 + m_L
\end{bmatrix}
\end{equation}

حيث يمثل العنصر $J_{3z}$ عزم القصور الذاتي الإجمالي الفعال والمصحح حول المحور العمودي، ويُحسب بالاعتماد على انزياح الاسطوانة وانزياح الحمل عن المركز المشترك :
\begin{equation}
J_{3z} = \frac{1}{3}m_3 L_3^2 + 2m_3 h_3^2 + 2m_L h_L^2 - \frac{1}{2}m_3 r^2
\end{equation}

حيث أن $h_3 = \left| \frac{L_3}{2} - z_{c3}^* \right|$ يمثل انزياح مركز ثقل الوصلة الأصلية، و $h_L = \left| L_3 - z_{c3}^* \right|$ يمثل انزياح مركز ثقل الحمل، و $r$ هو نصف قطر اسطوانة الوصلة الثالثة.
\section{اشتقاق مصفوفات النموذج الديناميكي }
\label{sec:lagrange_derivation}

بناءً على صياغة لاغرانج-أويلر التي تم تقديمها في القسم الأول، يتطلب حساب القوى والعزوم المحركة تفكيك تابع لاغرانج وحساب مشتقاته الجزئية بالنسبة لمتغيرات المفاصل وسرعاتها.  يمكن كتابة العزم الكلي الفعال المطبق على المفصل $(i)$ قبل دمج المحركات على النحو التالي:

\begin{equation}
\tau_i = \sum_{j=1}^{n} D_{ij}\ddot{q}_j + \sum_{j=1}^{n} \sum_{k=1}^{n} D_{ijk}\dot{q}_j\dot{q}_k + D_i
\label{eq:full_analytical_lagrange}
\end{equation}


\subsection{عناصر مصفوفة العطالة $(D_{ij})$}
تمثل المعاملات $(D_{ij})$ عناصر مصفوفة القصور الذاتي الكلية للروبوت $M(q)$ بأبعاد $3 \times 3$. تُعبر هذه العناصر عن عطالة المحاور والتداخل القصور الذاتي المتبادل بين المفاصل أثناء التسارع. تُحسب هذه المعاملات تحليلياً بالاعتماد على أثر مصفوفة التحويل المتكاملة (\LR{Trace}) ومصفوفات الاشتقاق الحركي $(U_{pi})$ وفق القانون التالي:

\begin{equation}
D_{ij} = \sum_{p=\max(i,j)}^{n} \text{Trace} \left( U_{pj} J_p U_{pi}^T \right)
\label{eq:inertia_element_law}
\end{equation}

حيث يمثل $J_p$ مصفوفة القصور الذاتي المتكاملة للوصلة $(p)$ والتي تم تحديدها في القسم السابق، بينما تمثل المصفوفة $U_{pi}$ المشتق الجزئي لمصفوفة التحويل الكلية بالنسبة للمتغير المفصلي $q_i$، والتي تعكس مدى تأثر تموضع الوصلة بتغير زاوية أو إزاحة المفصل.

\subsection{معاملات قوى كوريوليس والطرد المركزي $(D_{ijk})$}
تُعبر المعاملات ثلاثية الدليل $(D_{ijk})$ عن تأثيرات قوى كوريوليس (عند $j \neq k$) وقوى الطرد المركزي (عند $j = k$) الناتجة عن التداخل الحركي للوصلات الدورانية أثناء الدوران السريع. تُحسب هذه المعاملات قياسياً عبر العلاقة:

\begin{equation}
D_{ijk} = \sum_{p=\max(i,j,k)}^{n} \text{Trace} \left( U_{pjk} J_p U_{pi}^T \right)
\label{eq:coriolis_element_law}
\end{equation}

وكما تم إقراره في الفرضيات الأساسية لهذا الفصل، ونظراً لتشغيل الروبوت ضمن مجالات سرعة منخفضة إلى متوسطة، فقد تم إهمال هذه الحدود بالكامل في هذا المشروع، مما يعني هندسياً أن:
\begin{equation}
\sum_{j=1}^{n} \sum_{k=1}^{n} D_{ijk}\dot{q}_j\dot{q}_k = 0
\end{equation}

\subsection{عناصر شعاع الجاذبية الأرضية $(D_i)$}
تمثل المعاملات $(D_i)$ عناصر شعاع العزوم الناتجة عن الجاذبية الأرضية $G(q)$ المؤثرة على مفاصل الذراع. بالاعتماد على متجهات مراكز الثقل المحددة سابقاً $r_p$ لكل وصلة، واعتبار الأرض هي المرجعية الصفرية للارتفاع، تم استنتاج عناصر شعاع الجاذبية تحليلياً وفق القانون التالي:

\begin{equation}
D_i = \sum_{p=i}^{n} -m_p g^T U_{pi} r_p
\label{eq:gravity_element_law}
\end{equation}

حيث أن $g = [0, 0, -g]^T$ هو شعاع تسارع الجاذبية الأرضية، و $m_p$ هي كتلة الوصلة (مع أخذ الكتلة المركبة المتغيرة $m_3^*$ للوصلة الثالثة بعين الاعتبار).

\subsection{النمذجة المصفوفية المستنتجة للروبوت}
بتعويض القوانين التحليلية السابقة وتجميع العناصر، يختزل النموذج الديناميكي الصرف لذراع الروبوت (قبل ربطه بالجهود الكهربائية للمحركات) إلى البنية المصفوفية المدمجة التالية:

\begin{equation}
\begin{bmatrix} \tau_1 \\ \tau_2 \\ \tau_3 \end{bmatrix}_{robot} = 
\begin{bmatrix} 
D_{11} & D_{12} & D_{13} \\ 
D_{21} & D_{22} & D_{23} \\ 
D_{31} & D_{32} & D_{33} 
\end{bmatrix}
\begin{bmatrix} \ddot{\theta}_1 \\ \ddot{\theta}_2 \\ \ddot{d}_3 \end{bmatrix} + 
\begin{bmatrix} D_1 \\ D_2 \\ D_3 \end{bmatrix}
\label{eq:robot_matrix_form}
\end{equation}

تُشكل هذه المنظومة الرياضية المقيدة الرابط الفيزيائي الأساسي بين حركة ذراع الروبوت الميكانيكية والعزوم الخارجية الفعالة المطلوبة لتنفيذ المناورة التشغيلية.
\section{دمج المشغلات بالنموذج الديناميكي}
\label{sec:actuator_dynamics}

إن النموذج الديناميكي المستنتج في المعادلة (\ref{eq:robot_matrix_form}) يصف العلاقة بين حركة الروبوت والعزوم الميكانيكية المطبقة على مفاصله. ومع ذلك، في التطبيقات الصناعية الفعلية، يتم تشغيل المفاصل باستخدام محركات كهربائية مقترنة بعلب سرعة (\LR{Gearboxes}). وبالتالي، فإن إشارة التحكم الحقيقية هي الجهد الكهربائي $(u)$ المطبق على المحركات. للحصول على تمثيل دقيق للنظام بالكامل، لا بد من دمج ديناميكية هذه المحركات مع ديناميكية الوصلات الميكانيكية.

\subsection{النموذج الرياضي لمحرك التيار المستمر (\LR{DC Motor Model})}
بإهمال التحريضية الكهربائية للمحرك (\LR{Armature Inductance}) لكون الثابت الزمني الكهربائي أصغر بكثير من الثابت الزمني الميكانيكي، يمكن كتابة معادلة العزم الصافي الذي يولده المحرك وينقله إلى مفصل الروبوت بالصيغة التالية:

\begin{equation}
\tau_{act} = \left(\frac{K N}{R_a}\right) u - \left(N^2 J_m\right) \ddot{q} - N^2 \left(f_m + \frac{K^2}{R_a}\right) \dot{q}
\label{eq:motor_torque}
\end{equation}

حيث:
\begin{itemize}
    \item $(u)$: جهد التحكم المطبق على المحرك.
    \item $(K)$: ثابت العزم والقوة المحركة الكهربائية العكسية (\LR{Back-EMF Constant}).
    \item $(N)$: نسبة التخفيض لعلبة السرعة (\LR{Gear Ratio}).
    \item $(R_a)$: المقاومة الكهربائية لملفات المحرك.
    \item $(J_m)$: عزم القصور الذاتي العائد لدوار المحرك.
    \item $(f_m)$: معامل الاحتكاك اللزج الخاص بالمحرك.
\end{itemize}

ولتبسيط التدوين الرياضي تمهيداً لدمج المعادلات، تم تجميع الثوابت الفيزيائية للمحرك في ثلاثة معاملات رئيسية $(A, B, c)$ كما يلي:
\begin{equation}
A = \frac{K N}{R_a} \quad , \quad B = N^2 J_m \quad , \quad c = N^2 \left(f_m + \frac{K^2}{R_a}\right)
\end{equation}
وبالتالي تصبح معادلة المحرك:
\begin{equation}
\tau_{act} = A u - B \ddot{q} - c \dot{q}
\label{eq:motor_simplified}
\end{equation}

\subsection{النموذج الديناميكي الشامل والمدمج (\LR{Comprehensive Dynamic Model})}
لإيجاد توازن العزوم الكلي، نقوم بمساواة العزم المحرك المولد من الآلية (\ref{eq:motor_simplified}) مع العزم المقاوم الذي يتطلبه الروبوت للحركة (\ref{eq:robot_matrix_form})، مع إضافة قوى الاحتكاك اللزج الميكانيكي للمفاصل الممثلة بالمعامل $(f)$. تصبح معادلة التوازن لكل مفصل $(i)$:

\begin{equation}
D_{i1}\ddot{q}_1 + D_{i2}\ddot{q}_2 + D_{i3}\ddot{q}_3 + f_i \dot{q}_i + D_i = A_i u_i - B_i \ddot{q}_i - c_i \dot{q}_i
\end{equation}

بنقل جميع الحدود المضروبة بالتسارع $(\ddot{q})$ والسرعة $(\dot{q})$ إلى طرف واحد، نلاحظ أن عطالة المحرك $(B)$ تُضاف مباشرة إلى القطر الرئيسي لمصفوفة القصور الذاتي للروبوت $(D_{ii})$، وأن احتكاك المحرك $(c)$ يُضاف إلى احتكاك المفصل $(f)$. لتسهيل الصياغة، نُعرّف المعاملات المدمجة التالية:
\begin{equation}
\bar{M}_{ii} = D_{ii} + B_i \quad , \quad R_i = f_i + c_i
\end{equation}

وبإعادة كتابة النظام بأكمله بصيغة شعاعية مصفوفية مدمجة، نحصل على:
\begin{equation}
\bar{M}(q) \begin{bmatrix} \ddot{\theta}_1 \\ \ddot{\theta}_2 \\ \ddot{d}_3 \end{bmatrix} = 
\begin{bmatrix} A_1 u_1 - R_1 \dot{\theta}_1 - D_1 \\ A_2 u_2 - R_2 \dot{\theta}_2 - D_2 \\ A_3 u_3 - R_3 \dot{d}_3 - D_3 \end{bmatrix}
\end{equation}

حيث تمثل $\bar{M}(q)$ مصفوفة العطالة الكلية للنظام (الروبوت + المحركات). وللحصول على شعاع التسارع كخرج صريح جاهز للاستخدام في محاكاة \LR{Simulink} وتصميم خوارزميات التحكم المتقدمة، نضرب طرفي المعادلة بمقلوب مصفوفة العطالة $\bar{M}^{-1}$:

\begin{equation}
\begin{bmatrix} \ddot{\theta}_1 \\ \ddot{\theta}_2 \\ \ddot{d}_3 \end{bmatrix} = \bar{M}^{-1} 
\begin{bmatrix} A_1 u_1 - R_1 \dot{\theta}_1 - D_1 \\ A_2 u_2 - R_2 \dot{\theta}_2 - D_2 \\ A_3 u_3 - R_3 \dot{d}_3 - D_3 \end{bmatrix}
\label{eq:final_state_space}
\end{equation}

تُعد المعادلة (\ref{eq:final_state_space}) تتويجاً لجهود النمذجة الديناميكية في هذا الفصل. فهذا النموذج الشامل، الذي يأخذ بعين الاعتبار تغيرات كتلة الحمل، وانزياح مراكز الثقل، والقصور الذاتي اللحظي، وديناميكية المحركات الكهربائية، يشكل بيئة الاختبار (\LR{Plant}) الدقيقة التي ستخضع لاستراتيجيات التحكم المتين والتكيفي في الفصول القادمة.